\subsection{Introduction: Beyond the Black Box \quad / \quad \textthai{บทนำ: ก้าวข้ามกล่องดำ}}
\begin{paracol}{2}
The primary objective of classical physics is to distill complex observations into simple, elegant mathematical expressions. However, modern Deep Learning (DL) often moves in the opposite direction, creating "Black Box"\index{Black Box@กล่องดำ (Black Box)} models with millions of parameters that are opaque to human intuition. While a neural network can predict a trajectory with high precision, it does not provide us with the \textit{fundamental law}—the force or energy relation—governing that motion.

This chapter explores \textbf{Symbolic Regression (SR)}\index{Symbolic Regression@การถดถอยเชิงสัญลักษณ์ (Symbolic Regression)}, a paradigm shift in AI-driven physics. Unlike standard regression, which optimizes coefficients for a fixed formula, SR searches the space of all possible mathematical expressions. It uses evolutionary algorithms to "evolve" equations that are both accurate and simple.

We will focus on \textbf{Explainable AI (xAI)}\index{xAI@xAI (xAI)} through discovery. By the end of this chapter, we will demonstrate how AI can rediscover Kepler's Third Law from raw orbital data, bridging the gap between data-driven machine learning and symbolic theoretical physics.

\switchcolumn

\begin{thai}
วัตถุประสงค์หลักของฟิสิกส์ดั้งเดิมคือการสรุปการสังเกตที่ซับซ้อนให้เป็นนิพจน์ทางคณิตศาสตร์ที่เรียบง่ายและสง่างาม อย่างไรก็ตาม ดีปเลิร์นนิงสมัยใหม่มักจะดำเนินไปในทิศทางตรงกันข้าม โดยสร้างโมเดล "กล่องดำ (Black Box)" ที่มีพารามิเตอร์หลายล้านตัวซึ่งยากต่อการเข้าใจด้วยสัญชาตญาณของมนุษย์ แม้ว่าโครงข่ายประสาทเทียมจะสามารถทำนายวิถีการเคลื่อนที่ได้อย่างแม่นยำสูง แต่มันไม่ได้ให้ "กฎพื้นฐาน" แก่เรา เช่น ความสัมพันธ์ของแรงหรือพลังงานที่ควบคุมการเคลื่อนที่นั้น

บทนี้จะสำรวจ \textbf{การถดถอยเชิงสัญลักษณ์ (Symbolic Regression - SR)} ซึ่งเป็นการเปลี่ยนกระบวนทัศน์ในฟิสิกส์ที่ขับเคลื่อนด้วย AI ต่างจากการถดถอยมาตรฐานที่จะปรับค่าสัมประสิทธิ์สำหรับสูตรที่กำหนดไว้ตายตัว SR จะค้นหาในปริภูมิของนิพจน์ทางคณิตศาสตร์ที่เป็นไปได้ทั้งหมด โดยใช้อัลกอริทึมเชิงวิวัฒนาการเพื่อ "วิวัฒน์" สมการที่มีทั้งความแม่นยำและเรียบง่าย

เราจะมุ่งเน้นไปที่ \textbf{AI ที่อธิบายได้ (Explainable AI - xAI)} ผ่านการค้นพบ เมื่อจบหลักสูตรนี้ เราจะแสดงให้เห็นว่า AI สามารถค้นพบกฎข้อที่สามของเคปเลอร์ขึ้นมาใหม่จากข้อมูลทางวงโคจรดิบได้อย่างไร ซึ่งเป็นการเชื่อมช่องว่างระหว่างการเรียนรู้ของเครื่องที่ขับเคลื่อนด้วยข้อมูลและฟิสิกส์เชิงทฤษฎีสัญลักษณ์
\end{thai}
\end{paracol}
